\documentclass{article}

\usepackage[utf8]{inputenc}     % pdfLaTeX
\usepackage[T1]{fontenc}
\usepackage[magyar]{babel}

\usepackage{amsthm}
\usepackage{amsmath}
\usepackage{amsfonts}
\usepackage{tikz}
\usetikzlibrary{arrows.meta, positioning}
\usepackage{graphicx}

\title{Combinatorial problems}
\author{Robin Rovenszky}
\date{Last updated: May 11, 2026}

\begin{document}

\maketitle

\section{Problem 1.}
We are given numbers $n_i \in [1, 7]$ with $n_i \in \mathbb{N^+}$. We write down the digits $n_1$ through $n_7$ in some order, or in other words, we permute them. Then we count how many times the following relation holds.
\begin{equation}
n_i = n_{i+1} - 1.
\end{equation}
Then we define $N(n, k)$ to be the number of permutations where relation (1) is satisfied exactly $k$ times, with $max(n_i) = n$.
Since for $n = 7$, the number of permutations ($7! = 5040$) is already much larger than to be ideal for case-by-case analysis, we seek a recursive formula for the value of $N(n, k)$.

\end{document}